
\chapter{Getting Started}
\label{chap:getting-started}

This chapter explains how to obtain FEMaster, compile it from source, and
perform a first finite-element analysis. FEMaster is a command-line
application that reads a keyword-based input file, performs the requested
calculations, and writes the resulting data to output files. No graphical
interface is required, allowing the solver to be used both interactively
and as part of automated simulation workflows.


\section{Installation}

FEMaster is available for Linux, Windows, and macOS. The source code,
precompiled releases, and additional project information are hosted on
GitHub at \url{https://github.com/Luecx/FEMaster}.

The simplest way to get started is to download a precompiled version
from \url{https://github.com/Luecx/FEMaster/releases}. Where available,
release archives contain executables for the corresponding operating
system and architecture. After extracting the archive, the executable
can be launched directly from a terminal. Depending on the selected
release, additional runtime libraries may be required.

Alternatively, FEMaster can be compiled from source. This allows the
solver to be configured for the available hardware, including optional
multithreaded CPU and GPU acceleration, and provides access to the
latest development version. The source code can be obtained through
Git:

\begin{shellcmd}
git clone https://github.com/Luecx/FEMaster.git
cd FEMaster
\end{shellcmd}


\section{Building FEMaster}

FEMaster uses CMake to manage its compilation and external dependencies.
The repository provides predefined build configurations, called
\emph{presets}, for the supported operating systems and numerical
backends. A preset selects the compiler configuration, optimization
settings, and optional libraries required for a particular build.

From the repository root, configuring and compiling FEMaster requires
only two commands. For example, a basic Linux build can be created
using:

\begin{shellcmd}
cmake --preset linux
cmake --build --preset linux -j8
\end{shellcmd}

The name \val{linux} is the selected preset and can be replaced by any
of the available presets listed below. The first command configures
the project and prepares its dependencies, while the second compiles
the executable. The optional argument \val{-j8} specifies eight
parallel build jobs. Subsequent builds generally require only the
second command, unless the build configuration changes.

Compilation requires CMake 3.24 or newer, Git, and a compiler supporting
C++17. The presets using Ninja additionally require Ninja to be
installed. Standard dependencies, including Eigen, Spectra, argparse,
and LZ4, are downloaded automatically by CMake during configuration.
Optional numerical libraries such as Intel oneMKL, NVIDIA CUDA, and
cuDSS must be installed separately when their corresponding backends
are enabled. On macOS, Apple Accelerate is provided by the system SDK.


\subsection{Available build presets}

The predefined presets provide different combinations of CPU
parallelization, sparse solvers, and GPU acceleration. All regular
presets produce optimized Release builds, with the exception of the
profiling configuration. The available presets are summarized below.

\begin{longtable}{@{}p{0.36\linewidth}p{0.56\linewidth}@{}}
\toprule
\textbf{Preset} & \textbf{Configuration} \\
\midrule
\endhead

\multicolumn{2}{@{}l}{\textbf{Linux}} \\
\midrule
\val{linux} & Standard CPU build. \\
\val{linux-omp} & CPU build with OpenMP parallelization. \\
\val{linux-profile} & CPU build with profiling and debugging symbols. \\
\val{linux-mkl} & CPU build with sequential Intel oneMKL. \\
\val{linux-mkl-omp} & Threaded Intel oneMKL with OpenMP. \\
\val{linux-cuda} & NVIDIA CUDA acceleration. \\
\val{linux-cuda-cudss} & CUDA with NVIDIA cuDSS sparse direct solving. \\
\val{linux-cuda-mkl} & CUDA with sequential Intel oneMKL. \\
\val{linux-cuda-mkl-omp} & CUDA with threaded Intel oneMKL and OpenMP. \\
\val{linux-full} & CUDA, cuDSS, threaded Intel oneMKL, and OpenMP. \\

\midrule
\multicolumn{2}{@{}l}{\textbf{Windows}} \\
\midrule
\val{windows-msvc} & Standard CPU build using MSVC. \\
\val{windows-msvc-cuda} & MSVC build with CUDA acceleration. \\
\val{windows-msvc-cuda-cudss} & MSVC build with CUDA and cuDSS. \\
\val{windows-intel} & Standard CPU build using IntelLLVM. \\
\val{windows-intel-mkl} & IntelLLVM with sequential Intel oneMKL. \\
\val{windows-intel-mkl-omp} & IntelLLVM with threaded Intel oneMKL and OpenMP. \\

\midrule
\multicolumn{2}{@{}l}{\textbf{macOS}} \\
\midrule
\val{macos} & Standard CPU build. \\
\val{macos-omp} & CPU build with OpenMP parallelization. \\
\val{macos-accelerate} & CPU build with Apple Accelerate. \\
\val{macos-omp-accelerate} & OpenMP with Apple Accelerate. \\

\bottomrule
\end{longtable}

The available configurations depend on the installed toolchains and
hardware. Windows MSVC presets use Visual Studio, while the IntelLLVM
presets require the Intel compiler toolchain. CUDA configurations
require a compatible NVIDIA GPU and CUDA installation, and cuDSS
configurations additionally require the cuDSS library. Apple
Accelerate is available only on macOS.

CMake places the compiled executable in the build directory
associated with the selected preset, typically
\file{build/<preset>/bin}. Multi-configuration generators, such as
the Visual Studio presets, place the Release executable in
\file{build/<preset>/bin/Release}. The executable is named
\file{FEMaster} on Linux and macOS and \file{FEMaster.exe} on Windows.

The exact preset definitions and additional configuration options
can be found in \file{CMakePresets.json} and \file{CMakeLists.txt}.


\section{Running FEMaster}

FEMaster is executed by passing the path of an input deck to the
command-line application. Input files conventionally use the
\file{.inp} extension and contain the model definition together
with the requested analysis procedures. Both supported Abaqus
keywords and FEMaster-specific extensions are interpreted by the
same parser, so no explicit input-format selection is required.

If the executable is available on the system path, a model can
be executed using:

\begin{shellcmd}
FEMaster example.inp
\end{shellcmd}

Alternatively, the executable can be called using its location
in the build directory. For example, from the repository root
after compiling with the \val{linux} preset:

\begin{shellcmd}
./build/linux/bin/FEMaster example.inp
\end{shellcmd}

Several optional command-line arguments control execution and
result storage. The most commonly used options are:

\begin{longtable}{@{}p{0.31\linewidth}p{0.61\linewidth}@{}}
\toprule
\textbf{Argument} & \textbf{Description} \\
\midrule
\endhead
\key{--ncpus N} & Set the number of CPU threads available to the solver. \\
\key{--output FILE} & Specify an alternative output file base name or path. \\
\key{--output-format} & Select one or more result formats: \val{res}, \val{frd}, or \val{femr}. \\
\bottomrule
\end{longtable}

For example, the following command limits execution to eight CPU
threads and requests FRD and FEMR output:

\begin{shellcmd}
FEMaster example.inp --ncpus 8 --output-format frd femr
\end{shellcmd}

The selected numerical backend is configured through the input
deck and must be supported by the compiled executable. The
complete command-line interface can be inspected using
\key{--help}.


\section{First example: linear static truss}
\label{sec:first-example}

As a first calculation, consider a straight truss element subjected
to an axial tensile force. The model consists of two nodes connected
by a single element with a length of 1000 mm and a cross-sectional
area of 100 mm\(^2\). The material is linear elastic steel with
Young's modulus of 210000 MPa. One end is fixed, while the opposite
end is allowed to move only in the axial direction and is subjected
to a force of 1000 N.

Save the following input deck as \file{example.inp}:

\begin{syntax}
*HEADING
First FEMaster example

*NODE
1, 0., 0., 0.
2, 1000., 0., 0.

*ELEMENT, TYPE=T3, ELSET=BAR
1, 1, 2

*MATERIAL, NAME=STEEL
*ELASTIC
210000., 0.3

*TRUSS SECTION, ELSET=BAR, MATERIAL=STEEL
100.

*STEP, NAME=AXIAL_TENSION
*STATIC
*BOUNDARY
1, 1, 3, 0.
2, 2, 3, 0.
*CLOAD
2, 1, 1000.
*END STEP
\end{syntax}

The example combines FEMaster-specific element and section
definitions with Abaqus-style analysis keywords. Both are
supported within the unified input language. To perform the
analysis, execute:

\begin{shellcmd}
FEMaster example.inp
\end{shellcmd}

FEMaster reads the model, assembles the finite-element system,
solves the static equilibrium equations, and writes the calculated
results to output files. The expected axial displacement at the
loaded node is approximately \(0.047619\) mm, providing a simple
reference value for checking the calculation.


\section{Result files}

By default, FEMaster writes the results of an analysis using the
input filename as the output base name. Running \file{example.inp}
therefore produces \file{example.res} and \file{example.frd}.
An additional binary result format, \file{.femr}, is available
through the \key{--output-format} option.

The \file{.res} format stores numerical result Fields in a
text-based representation, while \file{.frd} provides results in
the CalculiX/CGX format for visualization and postprocessing.
The \file{.femr} format provides structured binary storage with
compression for more efficient handling of larger result datasets.
The available result quantities, output controls, and file
conventions are described in Section~\ref{sec:output-results}.


\section{Built-in documentation}

In addition to this manual, FEMaster provides a command-line
documentation mode for inspecting the keywords and input
variants recognized by the executable. This can be useful when
working with an unfamiliar command or checking which syntax
forms are supported by a particular build.

Documentation mode is activated using \key{--document}, followed
by an action specifying the requested information:

\begin{shellcmd}
FEMaster --document --doc-list
FEMaster --document --doc-show CLOAD
FEMaster --document --doc-variants ELASTIC
FEMaster --document --doc-search SECTION
\end{shellcmd}

These commands list available keywords, display the definition
of a particular command, inspect its accepted data variants,
or search the command registry. Documentation mode does not
require an input deck and does not perform a finite-element
analysis. For detailed explanations, modelling conventions,
and practical examples, the Keyword Reference remains the
primary source of information.
